\documentclass[11pt]{article}
\usepackage[T1]{fontenc}
\usepackage[utf8]{inputenc}
\usepackage[margin=26mm]{geometry}
\usepackage{amsmath,amssymb,amsthm,booktabs}
\usepackage[hidelinks]{hyperref}
\usepackage{microtype}
\newtheorem{theorem}{Theorem}
\newcommand{\gin}{\operatorname{gin}_{\mathrm{revlex}}}
\newcommand{\HS}{\operatorname{HS}}
\newcommand{\GL}{\operatorname{GL}}
\title{A seven-hyperplane counterexample to a generic-initial degree bound}
\author{Alec Kriebel\\
\small Independent researcher\\
\small ORCID: \href{https://orcid.org/0009-0001-9320-500X}{0009-0001-9320-500X}}
\date{Research note, 15 October 2026}
\begin{document}
\maketitle
\begin{abstract}
For the characteristic-zero arrangement defined by
$Q=xy(x-y)(x-2y)zw(x+z+w)$, the degree-reverse-lexicographic
generic initial ideal of the ambient Jacobian ideal has pure-second-variable
cutoff $8$ and a degree-$7$ minimal generator whose largest variable is the
third. This explicitly contradicts the second-variable bound of
Bigatti--Palezzato--Torielli, also stated in OWR Report 5/2021.
We prove the generic sectional dimensions on a nonempty open set of compatible
flags, using scalar polar syzygies and Hilbert--Burch, and give exact rational
checks. The arrangement is a classical parallel connection; the methods and
the valid three-variable arrangement theorem are prior work. Historical
novelty and independent discovery are not established. This unrefereed note
was developed with extensive AI assistance and has not undergone conventional
human peer review.
\end{abstract}

\section{Statement and conventions}
Let $K$ be any field of characteristic zero. Use the order
$x_1>x_2>x_3>x_4$ for degree reverse lexicographic generic initial ideals.
For an arrangement $\mathcal A$ with Jacobian ideal $J$, put
\[
 B=\gin(J),\qquad p=\min\{d\geq 0:x_2^d\in B\}.
\]
Conjecture 5.7 of Bigatti--Palezzato--Torielli \cite{BPT} asks that a
minimal generator of $B$ involving $x_3$ have degree at least $p$.
Their notation is $d_0=\min\{d:x_2^{d+1}\in B\}$, so $p=d_0+1$.
The same second-variable statement appears as Conjecture 7 on printed
page 234 of OWR Report 5/2021 \cite{OWR}, and as Conjecture 7.9 in
\cite{PT2020,PT2022}. A first-variable exponent in an imported problem
record was a transcription error; correcting it is source normalization,
not a mathematical contribution.

\begin{theorem}\label{thm:main}
In $S=K[x,y,z,w]$, let
\[
 Q=xy(x-y)(x-2y)zw(x+z+w),\qquad
 J=(Q_x,Q_y,Q_z,Q_w).
\]
For $B=\gin(J)\subset K[x_1,x_2,x_3,x_4]$, one has $p=8$,
and $B$ has a degree-$7$ minimal generator in $K[x_1,x_2,x_3]$
divisible by $x_3$. In particular its largest variable is $x_3$ and its
degree is strictly less than $p$.
\end{theorem}

The seven factors are distinct, and the coordinate factors make the central
arrangement essential. Euler's identity
$xQ_x+yQ_y+zQ_z+wQ_w=7Q$ shows that adjoining $Q$ does not change $J$.
The four polars have degree $6$.

The arrangement itself is a classical construction: it is the parallel
connection along the shared factor $x$ of the rank-two four-element pencil
$\{x,y,x-y,x-2y\}$ and the rank-three four-element circuit
$\{x,z,w,x+z+w\}$. Their matroids are $U_{2,4}$ and $U_{3,4}$.
See Falk--Proudfoot \cite[\S4, Theorem 4.1]{FP} for the realization of
parallel connections by identifying shared coordinates. No new-arrangement
claim is made here.

Marchesi--Palezzato--Torielli prove the conjecture for arrangements in $K^3$
\cite[Theorem 4.7]{MPT}. Our ternary section retains \emph{four ambient
polars}. It is not the three-generated Jacobian ideal obtained by first
restricting $Q$ and then differentiating in three variables. Theorem
\ref{thm:main} therefore respects that result.

\section{Polar relations on an open set of sections}
For an injective map $C:K^r\to K^4$, $r=2,3$, let $R_r$ be its polynomial
coordinate ring and write $X,Y,Z,W$ for the four coordinate restrictions.
Write $L_1,\ldots,L_7$ for the restricted factors, $a_i\in K^4$ for the
\emph{ambient} coefficient row of the $i$th factor, and
\[
 f_j=Q_j\circ C\quad (j=x,y,z,w),\qquad
 I_r=(f_x,f_y,f_z,f_w)\subset R_r.
\]
Require that $X,Y$ be independent and that all seven $L_i$ be nonzero and
pairwise nonproportional. For $r=3$ also require that $X,Z,W$ be independent.
These are open rank conditions. Nonemptiness for compatible sections is
proved below.

For a tuple $u=(u_x,u_y,u_z,u_w)$ of homogeneous linear forms in $R_r$,
ambient differentiation gives
\begin{equation}\label{eq:polar}
 \sum_j u_jf_j=\sum_{i=1}^7(a_i\cdot u)\prod_{h\ne i}L_h.
\end{equation}
Reduce a zero relation modulo $L_i$. The quotient $R_r/(L_i)$ is a domain,
and every other factor has nonzero image. Thus $L_i\mid a_i\cdot u$.
Since both have degree one, there is a scalar $\lambda_i$ with
$a_i\cdot u=\lambda_iL_i$. Conversely, these proportionalities turn
\eqref{eq:polar} into $(\sum_i\lambda_i)\prod_iL_i$. Consequently all
linear syzygies are characterized by
\begin{equation}\label{eq:scalar}
 a_i\cdot u=\lambda_iL_i\quad(1\leq i\leq7),\qquad
 \sum_i\lambda_i=0.
\end{equation}
The same reduction for a constant tuple forces $a_i\cdot u=0$ for every
$i$. The coordinate-factor rows imply $u=0$. Hence the four $f_j$ are
linearly independent and minimally generate $I_r$ in degree $6$.

This factor-divisibility calculation uses classical logarithmic-derivation
ideas. Tohaneanu \cite[\S2.1]{Tohaneanu} explains the linear case and
explicitly traces its decomposition mechanism to Orlik--Terao
\cite[Proposition 4.29(3)]{OT}. We apply that elementary mechanism to
restricted ambient polars; we do not assert a new general derivation theorem.

The first three factors $X,Y,X-Y$ force their scalars in
\eqref{eq:scalar} to coincide; $X-2Y$ has the same scalar, say $a$.
Call the remaining scalars $b,c,d$. The system becomes exactly
\begin{equation}\label{eq:reduced}
 \begin{gathered}
 u_x=aX,\quad u_y=aY,\quad u_z=bZ,\quad u_w=cW,\\
 (a-d)X+(b-d)Z+(c-d)W=0,\qquad 4a+b+c+d=0.
 \end{gathered}
\end{equation}

For $r=3$, independence of $X,Z,W$ gives $a=b=c=d$; the trace equation
is $7a=0$. Thus there is no linear syzygy. Writing
$H_r(k)=\dim_K(R_r/I_r)_k$, we obtain
\begin{equation}\label{eq:H3}
 H_3(5)=21,\qquad H_3(6)=28-4=24,\qquad H_3(7)=36-12=24.
\end{equation}

For $r=2$, the relation space of $X,Z,W$ is one-dimensional. Choose
$\rho_xX+\rho_zZ+\rho_wW=0$ nonzero. All solutions of the vector equation
in \eqref{eq:reduced} are
\[
 a=d+t\rho_x,\quad b=d+t\rho_z,\quad c=d+t\rho_w,
 \qquad 7d+t(4\rho_x+\rho_z+\rho_w)=0.
\]
The last equation has rank one, even if its coefficient of $t$ vanishes.
The map from its solution space to $u$ is injective: $u=0$ forces
$a=b=c=0$, and then $d(X+Z+W)=0$ forces $d=0$ and $t=0$.
There is therefore exactly one independent linear binary syzygy.

\section{Binary Hilbert function and genericity}
The ideal $I_2$ has height two. Indeed, work over an algebraic closure and
suppose all four binary polars vanish at a projective point. Restricted
Euler implies $\prod_iL_i=0$. Exactly one $L_i$ vanishes there, since
nonproportional binary linear forms have distinct zeros. The ambient
polar vector at that point is
$a_i\prod_{h\ne i}L_h$, a nonzero vector, which is a contradiction.
Thus there is no projective common zero, and $R_2/I_2$ is Artinian.
This conclusion descends to $K$ by extension of scalars.

The graded Hilbert--Burch theorem \cite{Burch} gives a minimal resolution
\[
 0\longrightarrow\bigoplus_{j=1}^3R_2(-6-e_j)
 \longrightarrow R_2(-6)^4\longrightarrow I_2\longrightarrow0,
 \qquad e_j\geq1.
\]
The finite-length quotient has Hilbert numerator
$N(t)=1-4t^6+\sum_jt^{6+e_j}$, which has a double zero at $1$.
The equation $N'(1)=0$ gives $\sum_je_j=6$. Exactly one linear syzygy
means exactly one $e_j$ is $1$; the other two are at least $2$.
Thus the coefficient degrees are $\{1,2,3\}$ and
\begin{equation}\label{eq:H2}
 \HS(R_2/I_2;t)=\frac{1-4t^6+t^7+t^8+t^9}{(1-t)^2},\qquad
 (H_2(6),H_2(7),H_2(8))=(3,1,0).
\end{equation}

For completeness the open conditions above hold on a compatible flag over
every field in scope. Take
\begin{equation}\label{eq:flag}
 (x,y,z,w)=(s,t,u,s+2t+3u),\qquad u=2s+3t.
\end{equation}
In the ternary section, $X,Y$ and $X,Z,W$ have the required independence.
The binary factor rows are
\[
 (1,0),\ (0,1),\ (1,-1),\ (1,-2),\ (2,3),\ (7,11),\ (10,14).
\]
They are nonzero with all pair determinants nonzero, as can be checked
directly. Pairwise distinctness also holds before the binary restriction.
A compatible full basis is the column basis of
\[
 C_0=\begin{pmatrix}
 1&0&0&0\\0&1&0&0\\2&3&1&0\\7&11&3&1
 \end{pmatrix},\qquad\det C_0=1.
\]
Its first two columns give the binary section and its first three span the
ternary section in \eqref{eq:flag}.

The conditions therefore define a nonempty open set of compatible flags,
whose lift to $\GL_4$ is nonempty open. The calculations
\eqref{eq:H3}--\eqref{eq:H2} hold throughout it. The variety $\GL_4$ is
irreducible, so this open set intersects the nonempty generic-initial open
set of Galligo's theorem \cite{Galligo}; see also \cite[Theorem 4.4]{BPT}.
Since a characteristic-zero field is infinite, the required nonempty opens
admit $K$-rational choices. Hence the displayed dimensions are the generic
sectional dimensions over every such $K$. The explicit flag proves
nonemptiness; favorable ranks at that flag alone would not prove genericity.

\section{The generic-initial generator}
The classical reverse-lexicographic section identity, for example
Conca \cite[Lemma 1.2]{Conca} and \cite[Theorem 6.2]{BPT}, identifies
the generic dimensions with
\begin{equation}\label{eq:section}
 H_r(k)=\dim_K\bigl(K[x_1,\ldots,x_r]/B_r\bigr)_k,
 \qquad B_r=B\cap K[x_1,\ldots,x_r].
\end{equation}
In characteristic zero, $B$ is strongly stable: if $x_j\mid m\in B$ and
$i<j$, then $x_im/x_j\in B$. Thus $x_2^k\in B$ implies that every binary
degree-$k$ monomial belongs to $B$. By \eqref{eq:H2}, $x_2^7\notin B$
and $x_2^8\in B$. Ideal closure rules out any earlier pure power, so $p=8$.

Put $A_3=K[x_1,x_2,x_3]/B_3$. The cokernel of
$x_3:(A_3)_{k-1}\to(A_3)_k$ is $(K[x_1,x_2]/B_2)_k$. Hence
\begin{equation}\label{eq:kernel}
 \dim_K\ker\bigl(x_3:(A_3)_{k-1}\to(A_3)_k\bigr)
 =H_3(k-1)-H_3(k)+H_2(k).
\end{equation}
This dimension is $21-24+3=0$ at $k=6$ and $24-24+1=1$ at $k=7$.
There are no elements of $B$ in degrees below $6$, since taking an initial
ideal preserves the graded Hilbert function. A degree-$6$ minimal generator
of $B_3$ divisible by $x_3$ would, after division by $x_3$, give a standard
monomial in the first kernel. Thus no generator of $B_3$ involving $x_3$
has degree at most $6$.

The second kernel is nonzero and is spanned by killed standard monomials.
Choose a degree-$6$ standard monomial $m$ with $x_3m\in B_3$, and let
$v$ be a minimal generator dividing $x_3m$. If $x_3\nmid v$, then $v\mid m$,
which contradicts standardness. Since $v$ involves $x_3$ and has degree at
most $7$, the preceding paragraph forces $\deg v=7$. It is also minimal in
$B$, because a divisor of a monomial in the first three variables cannot
involve $x_4$. This proves Theorem \ref{thm:main}.

Equation \eqref{eq:kernel} is the classical sectional-matrix triangle
calculation; compare \cite[Theorem 6.6]{BPT}. We use it inside $B_3$ and
then prove minimality in $B$ by divisibility. A row for the third truncation
does not detect arbitrary full-ring generators involving later variables.
No full generic initial ideal or particular exponent vector for $v$ is
needed.

\section{Exact verification}
The accompanying standard-library Python verification reconstructs $Q$ over
$\mathbb Q$, differentiates in four variables \emph{before} restriction,
and builds exact homogeneous multiplication matrices for \eqref{eq:flag}.
Their ranks are corroboration of the open-set proof:
\begin{center}
\begin{tabular}{ccrrr}
\toprule
Section variables & Total degree & Matrix shape & Rank & Quotient dimension\\
\midrule
3&6&$4\times28$&4&24\\
3&7&$12\times36$&12&24\\
2&6&$4\times7$&4&3\\
2&7&$8\times8$&7&1\\
2&8&$12\times9$&9&0\\
\bottomrule
\end{tabular}
\end{center}
Rows are polynomial multiples of the four polars and columns are degree-$k$
monomials. Exact nonzero minors and complete left-kernel bases certify both
rank bounds. The binary syzygy is
\[
 11s f_x+11t f_y-73(2s+3t)f_z+25(7s+11t)f_w=0.
\]
A constructive $4\times3$ Hilbert--Burch matrix with coefficient degrees
$(1,2,3)$ is also certified by all its signed maximal minors.
The files \texttt{reproduce\_exact.py} and \texttt{EXACT\_CERTIFICATE.json}
reconstruct and verify these data; \texttt{orthogonal\_checks.py} and
\texttt{ORTHOGONAL\_CERTIFICATE.json} add exact initial-ideal checks at
specified rational coordinates and excluded-case controls. See the package
README for commands requiring Python 3.11 or later. These finite coordinate
checks make no additional genericity or full-gin claim.

\section{Prior overlap, priority limits, and disclosure}
An applicable prior family theorem deserves explicit credit. For this
arrangement there is one codimension-two flat of multiplicity four, with
three hyperplanes outside the pencil. Migliore--Nagel
\cite[\S8.2, Theorem 8.11 and Corollary 8.12]{MN} give
\[
 0\longrightarrow S(-7)^3\oplus S(-9)
 \longrightarrow S(-6)^5\longrightarrow J^{\mathrm{top}}
 \longrightarrow0,
\]
where $J^{\mathrm{top}}$ is the intersection of the codimension-two primary
components of $J$. This is substantive prior information about this family.
It concerns a five-generated ideal, whereas $J$ has four minimal ambient
polars. Their \S3 also stresses the distinction between restricting an
ambient Jacobian ideal and taking the Jacobian ideal of a restricted
arrangement. The proof above concerns the successive sections of $J$.

Historical priority remains unresolved. A bounded literature audit completed
on 15 October 2026 did not identify an earlier complete predecessor in its
inspected corpus. In particular, the full text
of Ilardi's 2018 article \cite{Ilardi} was not obtained: publisher retrieval
failed and institutional access was restricted. Its abstract does not
exclude a relevant theorem or family. No assertion of firstness, independent
discovery, or present openness follows from the accessible literature.
The audit did not exhaust parameter families, related nine-factor generic
initial ideals, or historical code, and it did not authenticate the
candidate's first public disclosure timestamp.
The supported outcome here is an explicit counterexample with a checkable
proof. We make no claim of smallest hyperplane count, a full generic initial
ideal, or a positive-characteristic extension.

Extensive AI tools were used in solving, drafting, literature inspection,
exact-verification development, and adversarial checking. The artifact has
received multiple AI-assisted mathematical checks but no conventional human
peer review and is unrefereed. Those checks are not evidence of historical
priority or a substitute for such review. All cited construction and method
credits apply regardless of the eventual priority determination.

\begin{thebibliography}{99}\small
\bibitem{BPT}
A. M. Bigatti, E. Palezzato, and M. Torielli,
\emph{New characterizations of freeness for hyperplane arrangements},
J. Algebraic Combin. \textbf{51} (2020), 297--315.
\href{https://doi.org/10.1007/s10801-019-00876-9}{doi:10.1007/s10801-019-00876-9};
\href{https://arxiv.org/abs/1801.09868}{arXiv:1801.09868}.

\bibitem{OWR}
E. Palezzato (joint work with M. Torielli),
\emph{Hyperplane arrangements and $k$-Lefschetz properties}, pp. 233--235
in \emph{Logarithmic Vector Fields and Freeness of Divisors and Arrangements:
New perspectives and applications}, Oberwolfach Reports \textbf{18}
(2021), no. 1, Report 5/2021; Conjecture 7, p. 234.
\href{https://doi.org/10.4171/OWR/2021/5}{doi:10.4171/OWR/2021/5}.

\bibitem{PT2020}
E. Palezzato and M. Torielli,
\emph{Lefschetz properties and hyperplane arrangements},
J. Algebra \textbf{555} (2020), 289--304.
\href{https://doi.org/10.1016/j.jalgebra.2020.02.039}{doi:10.1016/j.jalgebra.2020.02.039}.

\bibitem{PT2022}
E. Palezzato and M. Torielli,
\emph{$k$-Lefschetz properties, sectional matrices and hyperplane arrangements},
J. Algebra \textbf{590} (2022), 215--233.
\href{https://doi.org/10.1016/j.jalgebra.2021.10.014}{doi:10.1016/j.jalgebra.2021.10.014}.

\bibitem{MPT}
S. Marchesi, E. Palezzato, and M. Torielli,
\emph{Lefschetz properties and the Jacobian algebra of $3$-dimensional
hyperplane arrangements}, Port. Math. \textbf{83} (2026), 1--18.
\href{https://doi.org/10.4171/PM/2155}{doi:10.4171/PM/2155}.

\bibitem{FP}
M. J. Falk and N. J. Proudfoot,
\emph{Parallel connections and bundles of arrangements},
Topology Appl. \textbf{118} (2002), 65--83.
\href{https://arxiv.org/abs/math/0002094}{arXiv:math/0002094}.

\bibitem{Tohaneanu}
S. O. Tohaneanu,
\emph{Projective duality of arrangements with quadratic logarithmic vector fields},
\href{https://arxiv.org/abs/1405.2122v2}{arXiv:1405.2122v2} (2015), \S2.1.

\bibitem{OT}
P. Orlik and H. Terao, \emph{Arrangements of Hyperplanes},
Grundlehren Math. Wiss. \textbf{300}, Springer, 1992.
\href{https://doi.org/10.1007/978-3-662-02772-1}{doi:10.1007/978-3-662-02772-1}.

\bibitem{Burch}
L. Burch, \emph{On ideals of finite homological dimension in local rings},
Proc. Cambridge Philos. Soc. \textbf{64} (1968), 941--948.
\href{https://doi.org/10.1017/S0305004100043620}{doi:10.1017/S0305004100043620}.

\bibitem{Galligo}
A. Galligo, \emph{\`A propos du th\'eor\`eme de pr\'eparation de Weierstrass},
in \emph{Fonctions de plusieurs variables complexes},
Lecture Notes in Math. \textbf{409}, Springer, 1974, 543--579.
\href{https://doi.org/10.1007/BFb0068121}{doi:10.1007/BFb0068121}.

\bibitem{Conca}
A. Conca, \emph{Reduction numbers and initial ideals},
Proc. Amer. Math. Soc. \textbf{131} (2003), 1015--1020.
\href{https://arxiv.org/abs/math/0210064}{arXiv:math/0210064}.

\bibitem{MN}
J. C. Migliore and U. Nagel,
\emph{Jacobian Ideals of Hyperplane Arrangements and their Graded Betti Numbers},
\href{https://arxiv.org/abs/2508.12113v1}{arXiv:2508.12113v1} (2025).

\bibitem{Ilardi}
G. Ilardi, \emph{Jacobian ideals, arrangements and the Lefschetz properties},
J. Algebra \textbf{508} (2018), 418--430.
\href{https://doi.org/10.1016/j.jalgebra.2018.04.029}{doi:10.1016/j.jalgebra.2018.04.029}.
Full text not obtained for the priority audit.
\end{thebibliography}
\end{document}
